% =====================================================================
%  INTRODUCCIÓN: párrafos continuos, sin subsecciones,
%  como pide la plantilla. Las citas nombran al autor dentro de la frase.
% =====================================================================
\chapter*{Introducción}
\addcontentsline{toc}{chapter}{INTRODUCCIÓN}

Un edificio es un entorno estructurado: sus pasillos, salones y escaleras siguen un
orden fijo. Aun así, quien entra por primera vez a uno grande construye un mapa mental con
información parcial y lo va corrigiendo mientras camina. Ese esfuerzo crece cuando el
espacio tiene pasillos repetidos, pocos puntos de referencia o señalización escasa, y
crece aún más al cambiar de piso, ya que la escalera interrumpe el recorrido. \textcite{gathmorad2024strategic} lo estudiaron con 149 participantes que
buscaban destinos en edificios virtuales de varios pisos, y encontraron que, entre pisos,
la eficiencia depende de lo que el edificio deja ver y de lo
que cada participante espera encontrar en él.

La robótica móvil colaborativa resuelve tareas parecidas, y el reparto
de trabajo entre varios robots tiene base teórica desde 2004, cuando
\textcite{gerkey2004taxonomy} formalizaron la asignación de tareas en sistemas
multirrobot. Sin embargo, casi todos esos sistemas operan en un solo piso. Cuando se
ha abordado el cambio de piso, la solución ha pasado por el hardware:
\textcite{kim2024indoor} combinan un mapa por piso con un grafo topológico y hacen que
el robot use el ascensor, mientras que \textcite{song2024elevator} automatizan ese uso
con dispositivos conectados que se instalan en el edificio. La otra opción es dotar al
robot de mecanismos para subir escaleras. En todos los casos aumentan los
requerimientos tecnológicos de la plataforma, y con ellos los costos, o se interviene la
infraestructura.

Este trabajo propone otro camino: en lugar de un robot que cambia de piso, dos
agentes, robots móviles terrestres, uno por piso. El usuario cambia de piso por su
cuenta y el agente del otro piso lo recibe en la escalera, de modo que el problema deja
de ser de locomoción y pasa a ser de coordinación. Ese traspaso del guiado entre dos
agentes, en un punto acordado de antemano, es el \emph{relevo}, y constituye el aporte
del proyecto.

La pregunta que organiza el trabajo es cómo asistir la orientación de usuarios dentro
de entornos interiores con múltiples pisos utilizando sistemas robóticos. Como respuesta se propone y se pone a prueba el relevo, para lo cual hay que demostrar que la
asistencia no se interrumpe cuando el usuario cambia de piso y medirlo con criterios
fijados antes de los experimentos.

Para implementarlo, el sistema se desarrolló con software libre y se cargó en dos agentes AWS DeepRacer, de Amazon Web Services. El
coordinador se comunica con cada agente por ROS~2, y el usuario pide el
servicio desde una página web, sin instalar nada y
sin salida a internet. Primero se evaluó en simulación, con treinta misiones sorteadas, y después con los agentes físicos, con diez,
entre los pisos 3 y 4 del edificio Catalayud de la sede principal de la Universidad Santo Tomás.
